\SourceLeaf{257}%
nec non per formulam \(r\cos\psi=\tfrac{1}{2}\mathord{.}\frac{\mathrm{d}rr}{\mathrm{d}q}\)
\[
\begin{aligned}
\relax[3]\quad r\cos\psi={}&q+\tfrac{2}{3}f^0ppq+\tfrac{1}{2}f'p^3q+(\tfrac{2}{5}f''-\tfrac{4}{45}f^0f^0)p^4q\quad\text{etc.}\\
&+\tfrac{3}{4}g^0ppqq+\tfrac{3}{5}g'p^3qq\\
&+(\tfrac{4}{5}h^0-\tfrac{14}{45}f^0f^0)ppq^3
\end{aligned}
\]
Hinc simul innotescit angulus \(\psi\). Perinde ad computum anguli \(\varphi\) concinnius evolvuntur series pro \(r\cos\varphi\) atque \(r\sin\varphi\), quibus inserviunt aequationes differentiales partiales
\[
\begin{aligned}
\frac{\mathrm{d}\mathord{.}r\cos\varphi}{\mathrm{d}p}&=n\cos\varphi\mathord{.}\sin\psi-r\sin\varphi\mathord{.}\frac{\mathrm{d}\varphi}{\mathrm{d}p}\\
\frac{\mathrm{d}\mathord{.}r\cos\varphi}{\mathrm{d}q}&=\cos\varphi\mathord{.}\cos\psi-r\sin\varphi\mathord{.}\frac{\mathrm{d}\varphi}{\mathrm{d}q}\\
\frac{\mathrm{d}\mathord{.}r\sin\varphi}{\mathrm{d}p}&=n\sin\varphi\mathord{.}\sin\psi+r\cos\varphi\mathord{.}\frac{\mathrm{d}\varphi}{\mathrm{d}p}\\
\frac{\mathrm{d}\mathord{.}r\sin\varphi}{\mathrm{d}q}&=\sin\varphi\mathord{.}\cos\psi+r\cos\varphi\mathord{.}\frac{\mathrm{d}\varphi}{\mathrm{d}q}\\
n\cos\psi\mathord{.}\frac{\mathrm{d}\varphi}{\mathrm{d}q}+\sin\psi\mathord{.}\frac{\mathrm{d}\varphi}{\mathrm{d}p}&=0
\end{aligned}
\]
quarum combinatio suppeditat
\[
\begin{aligned}
\frac{r\sin\psi}{n}\mathord{.}\frac{\mathrm{d}\mathord{.}r\cos\varphi}{\mathrm{d}p}+r\cos\psi\mathord{.}\frac{\mathrm{d}\mathord{.}r\cos\varphi}{\mathrm{d}q}&=r\cos\varphi\\
\frac{r\sin\psi}{n}\mathord{.}\frac{\mathrm{d}\mathord{.}r\sin\varphi}{\mathrm{d}p}+r\cos\psi\mathord{.}\frac{\mathrm{d}\mathord{.}r\sin\varphi}{\mathrm{d}q}&=r\sin\varphi
\end{aligned}
\]
Hinc facile evolvuntur series pro \(r\cos\varphi,r\sin\varphi\), quarum termini primi manifesto esse debent \(p\) et \(q\), puta
\[
\begin{aligned}
\relax[4]\quad r\cos\varphi={}&p+\tfrac{2}{3}f^0pqq+\tfrac{5}{12}f'ppqq+(\tfrac{3}{10}f''-\tfrac{8}{45}f^0f^0)p^3qq\quad\text{etc.}\\
&+\tfrac{1}{2}g^0pq^3+\tfrac{7}{20}g'ppq^3\\
&+(\tfrac{2}{5}h^0-\tfrac{7}{45}f^0f^0)pq^4
\end{aligned}
\]
\[
\begin{aligned}
\relax[5]\quad r\sin\varphi={}&q-\tfrac{1}{3}f^0ppq-\tfrac{1}{6}f'p^3q-(\tfrac{1}{10}f''-\tfrac{7}{90}f^0f^0)p^4q\quad\text{etc.}\\
&-\tfrac{1}{4}g^0ppqq-\tfrac{3}{20}g'p^3qq\\
&-(\tfrac{1}{5}h^0+\tfrac{13}{90}f^0f^0)ppq^3
\end{aligned}
\]
E combinatione aequationum [2], [3], [4], [5] derivari posset series pro \(rr\cos(\psi+\varphi)\), atque hinc, dividendo per seriem [1], series pro \(\cos(\psi+\varphi)\), a qua ad seriem pro ipso angulo \(\psi+\varphi\) descendere liceret. Elegantius tamen eadem obtinetur se%
